\documentclass[12pt,a4paper]{article}
\usepackage[utf8]{inputenc}
\usepackage[T5]{fontenc}
\usepackage[vietnamese]{babel}
\usepackage{amsmath, amssymb, amsfonts}
\usepackage[margin=1in]{geometry}
\usepackage{ulem}
\newcommand{\ul}[1]{\uline{#1}}

\begin{document}
\ul{Buổi 06 -- Ngày 23-08-2026 -- môn Đại số tuyến tính -- lớp MA003.F32.LT.CNTT}

\textbf{ĐỀ KIỂM TRA GIỮA KỲ III}

\textbf{NĂM HỌC 2025-2026 -- MÔN MA003}

(\emph{SV được phép sử dụng tài liệu khi làm bài})

Nộp bài trên Assignment của Elearning.

Làm bài và đặt tên file như sau:

\emph{\ul{Phương án 1}}:

Làm bài vào giấy $\rightarrow$ chụp hình $\rightarrow$ dán các hình bài làm vào 1 file word $\rightarrow$
lưu file theo định dạng .pdf theo quy ước ``MSSV\_Họ và tên.pdf'' $\rightarrow$ nộp
trên Assignment của Elearning.

\emph{\ul{Phương án 2}}:

Làm bài trên file Word/Excel/hoặc phần mềm soạn thảo văn bản khác $\rightarrow$ lưu
lại theo định dạng .pdf theo quy ước ``MSSV\_Họ và tên.pdf'' $\rightarrow$ nộp trên
Assignment của Elearning. Nhớ gõ công thức Toán bằng Microsoft Equation/
hoặc MathType.

\textbf{Chúc lớp mình vui khỏe, thành công} :) !

\emph{\textbf{\ul{Câu 1}}}: (\emph{3,0 điểm})

\begin{quote}
Cho các ma trận thực: \[A = \begin{pmatrix} 1 & 1 & 2 \\ 1 & 0 & m \\ -3 & -5 & -7 \end{pmatrix}\],
\[B = \begin{pmatrix} -1 & -4 & 2 \\ 5 & -9 & -1 \\ -3 & -6 & 0 \end{pmatrix}\],
\[C = \begin{pmatrix} 3 & -2 & \lambda \\ 1 & 1 & -2 \\ 4 & -1 & -2 \end{pmatrix}\],

với $m, \lambda$ là các tham số
thực.
\end{quote}

a/ Tìm ma trận: $2C^T B - 5A$.

b/ Tìm $\lambda$ để ma trận
$C$ có hạng nhỏ nhất.

c/ Với $m = 2$, tìm ma trận
$X$ thỏa
$XA = D + 3B$, với ma trận
\[D = \begin{pmatrix} -4 & -1 & 0 \\ -1 & 1 & -5 \\ 8 & -3 & 1 \end{pmatrix}\].

\emph{\textbf{\ul{Câu 2}}}: (\emph{3,0 điểm})

Giải và biện luận hệ phương trình tuyến tính sau trên trường số thực
$\mathbb{R}$:

\[\begin{cases} x - my + m^2 z = m \\ mx - m^2 y + mz = 1 \\ mx + y - m^2 z = 1 \end{cases}\], với
$m$ là tham số thực.

\emph{\textbf{\ul{Câu 3}}}: (\emph{2,0 điểm})

Trên không gian $\mathbb{R}^5$, cho
tập hợp \[W = \left\{ (x_1; x_2; x_3; x_4; x_5) \mid \begin{matrix} x_4 = 7x_1 - 4x_2 \\ x_1 + 3x_3 - 5x_2 = 0 \end{matrix} \right\}\].

Hỏi $W$ có phải là không
gian véc tơ con của $\mathbb{R}^5$
không? Vì sao?

\emph{\textbf{\ul{Câu 4}}}: (\emph{2,0 điểm})

Trên không gian $\mathbb{R}^4$, cho
tập hợp $S = \{u_1, u_2, u_3, u_4\}$, với
$u_1 = (0; 1; 3; 4)$,
$u_2 = (1; 0; 2; 3)$,
$u_3 = (-3; -2; 0; -5)$,
$u_4 = (4; 3; -5; 0)$.

a/ Hỏi tập hợp $S$ là độc
lập tuyến tính hay phụ thuộc tuyến tính? Vì sao?

\begin{quote}
b/ Hãy tìm một biểu diễn tuyến tính của véc tơ
$\lambda = 4\alpha - 3\beta$ theo các véc tơ
$u_1, u_2, u_3, u_4$, trong đó
$\alpha = (-10; -2; 15; -5)$ và
$\beta = (-2; 4; 10; -2)$.

-------- HẾT --------
\end{quote}

\textbf{Chương 3: KHÔNG GIAN VÉC TƠ (TIẾP THEO)}

\textbf{5/ \ul{KHÔNG GIAN CON SINH BỞI MỘT TẬP HỢP}}

Cho $(V, +, \bullet)$ là không gian
tuyến tính tổng quát trên trường số \emph{F},

và cho tập hợp $S = \{\alpha_1, \alpha_2, \dots, \alpha_m\}$ trên
\emph{V}.

Ta ký hiệu $\langle S \rangle = \{$tập hợp tất
cả các THTT có từ \emph{S}\}

$= \{ \alpha = c_1 \alpha_1 + c_2 \alpha_2 + \dots + c_m \alpha_m \mid c_1, c_2, \dots, c_m \in F; \alpha_1, \alpha_2, \dots, \alpha_m \in S \}$.

Ta có $\langle S \rangle \le V$, nghĩa là ta có
$\langle S \rangle$ là không gian véc tơ
con (không gian tuyến tính con) của \emph{V}, sinh bởi tập hợp \emph{S},
và thường viết là $\langle S \rangle$ is a
sub-vector space that is spanned by \emph{S}.

* \emph{\textbf{\ul{Ý nghĩa}}}:

Trên \emph{V} có thể có nhiều không gian con khác nhau của \emph{V} chứa
được tập hợp \emph{S}, nhưng
$\langle S \rangle$ là không gian con nhỏ
nhất, chứa được \emph{S}.

\emph{S}

\emph{V}

Ta còn gọi \emph{S} là hệ sinh (tập sinh) ra không gian tuyến tính
$\langle S \rangle$, và viết là \emph{S}
spans $\langle S \rangle$.

\textbf{6/ \ul{CƠ SỞ CỦA KHÔNG GIAN VÉC TƠ}}

Cơ sở (basis) của một không gian véc tơ (không gian tuyến tính) là một
hệ sinh (tập sinh), độc lập tuyến tính, của không gian tuyến tính đó.

Cơ sở của không gian tuyến tính không phải là duy nhất, nghĩa là có vô
số các cơ sở khác nhau của một không gian tuyến tính cho trước.

Nói cách khác, cơ sở của một gian véc tơ là tập hợp chứa những véc tơ cơ
bản nhất, cốt lõi nhất, để tạo ra một không gian véc tơ.

* \emph{\textbf{\ul{Ý nghĩa}}}:

Cho tập hợp $S = \{ \alpha_1, \alpha_2, \dots, \alpha_m \}$ là một cơ
sở của không gian véc tơ \emph{V}.

Lúc này, với mọi véc tơ
$\alpha \in V$, ta luôn tìm được bộ
số: $c_1, c_2, \dots, c_m$ sao cho

$\alpha = c_1 \alpha_1 + c_2 \alpha_2 + \dots + c_m \alpha_m$

nghĩa ta luôn biểu diễn được
$\alpha$ theo cơ sở \emph{S};
cho nên ta có thể phát sinh ra bất kỳ véc tơ nào trên \emph{V}, khi đã
biết cơ sở \emph{S}. Vì vậy, về bản chất là khi biết cơ sở của một không
gian tuyến tính nghĩa là ta đã kiểm soát được toàn bộ không gian tương
ứng.

\emph{\textbf{\ul{Lưu ý}}}:

Số lượng véc tơ trong cơ sở chính là số chiều của không gian tương ứng.

\textbf{* \ul{ĐỂ CHỨNG MINH MỘT TẬP HỢP CÓ PHẢI LÀ CƠ SỞ CỦA MỘT KHÔNG
GIAN VÉC TƠ HAY KHÔNG}}

Cho không gian tuyến tính
$(V, +, \bullet)$ tổng quát trên
\emph{F}.

Để chứng minh tập hợp $S = \{\alpha_1, \alpha_2, \dots, \alpha_m\}$
là một cơ sở của \emph{V}, ta làm như sau:

+ Ta lập luận \emph{S} là một hệ sinh ra \emph{V}.

+ Ta cần chứng tỏ \emph{S} là tập hợp ĐLTT.

Từ đó suy ra \emph{S} là một cơ sở của \emph{V}.

\ul{Ví dụ mẫu 8}: bài 3.12

Bài \textbf{3. 12.} Trong không gian $\mathbb{R}^3$ cho hệ véc tơ $\{a_1, a_2, a_3\}$ với
$a_1 = (1, 1, 1), \; a_2 = (1, 2, 3), \; a_3 = (2, 1, 4).$

Cho biết hệ $S = \{ a_1, a_2, a_3 \}$ có là một
cơ sở của $\mathbb{R}^3$ hay không?
Tại sao?

\ul{Giải}:

Ta có $\mathbb{R}^3$ là không gian 3
chiều trên $\mathbb{R}$, nghĩa là
cơ sở của $\mathbb{R}^3$ là một hệ
sinh chứa 3 véc tơ, độc lập tuyến tính.

Mà tập hợp $\{ a_1, a_2, a_3 \}$ có sẵn 3
véc tơ nên đây là một hệ sinh của
$\mathbb{R}^3$; cho nên để kiểm chứng
xem $\{ a_1, a_2, a_3 \}$ có là cơ sở của
$\mathbb{R}^3$ hay không thì ta chỉ
cần xét xem hệ $\{ a_1, a_2, a_3 \}$ có ĐLTT
hay không.

Ta lập ma trận

\begin{quote}
\[A_S = \begin{pmatrix} a_1 \\ a_2 \\ a_3 \end{pmatrix} = \begin{pmatrix} 1 & 1 & 1 \\ 1 & 2 & 3 \\ 2 & 1 & 4 \end{pmatrix}\]

Ta có

$\det(A_S) = 1.2.4 + 1.3.2 + 1.1.1 - [1.2.2 + 1.1.4 + 1.3.1] = 4 \ne 0$

Nên ta nói $\{ a_1, a_2, a_3 \}$ là ĐLTT;
cho nên hệ $\{ a_1, a_2, a_3 \}$ là một cơ
sở của $\mathbb{R}^3$.
\end{quote}

\textbf{* \ul{TÌM CƠ SỞ CHO KHÔNG GIAN KHI BIẾT TRƯỚC HỆ SINH (TẬP
SINH)}}

Cho $(V, +, \bullet)$ là một không gian
véc tơ (không gian tuyến tính) tổng quát trên \emph{F},

và cho tập hợp $S = \{\alpha_1, \alpha_2, \dots, \alpha_m\}$ trên
\emph{V}.

Gọi $W = \langle S \rangle$, nghĩa là \emph{S}
là hệ sinh (tập sinh) của \emph{W} (i.e. \emph{S} spans \emph{W})

Ta tìm cơ sở và xác định số chiều cho \emph{W} như sau:

\emph{\textbf{\ul{TH1}}}:
$m = 1$, nghĩa là
$S = \{ \alpha_1 \}$

Ta lập luận véc tơ $\alpha_1 \neq$véc
tơ $O$, suy ra
$S = \{ \alpha_1 \}$ vừa là hệ sinh, vừa là
cơ sở của \emph{W}, và số chiều của \emph{W} là:

$\dim_F W = 1$ (dim = dimension)

\emph{\textbf{\ul{TH2}}}:
$m = 2$, nghĩa là
$S = \{ \alpha_1, \alpha_2 \}$

Ta lập luận véc tơ $\alpha_1$
không tỷ lệ với $\alpha_2$ theo
tọa độ

(Ví dụ:
$\alpha_1 = (-3; 0; 1; 5; -6), \; \alpha_2 = (2; 1; -4; 8; 7)$,$\alpha_3 = (6; 0; -2; -10; 12)$.

Ta có: $\frac{-3}{2} \neq \frac{0}{1} \neq \frac{1}{-4} \neq \frac{5}{8} \neq \frac{-6}{7} \Rightarrow \alpha_1$ không tỷ lệ với
$\alpha_2$ theo tọa độ; nhưng
$\alpha_3$ tỷ lệ với
$\alpha_1$ do
$\alpha_3 = -2\alpha_1$).

Suy ra $S = \{\alpha_1, \alpha_2\}$ vừa là hệ sinh,
vừa là cơ sở của \emph{W}, và số chiều của \emph{W} là:

$\dim_F W = 2$.

\emph{\textbf{\ul{TH3}}}:
$m \ge 3$, nghĩa là
$S = \{\alpha_1, \alpha_2, \dots, \alpha_m\}$

Ta lập ma trận \[A_S = \begin{pmatrix} \alpha_1 \\ \alpha_2 \\ \vdots \\ \alpha_m \end{pmatrix}\]

\ul{TH3.1}: $A_S$ là một ma
trận vuông

Nếu $\det(A_S) \ne 0$ suy ra \emph{S} là
một hệ sinh, ĐLTT nên \emph{S} cũng là một cơ sở của \emph{W}, và

số chiều là $\dim_F W = m$.

Nếu $\det(A_S) = 0$ ta chuyển qua
TH3.2.

\ul{TH3.2}: $A_S$ là một ma
trận tùy ý, nghĩa là $A_S$
là một ma trận vuông hay không vuông đều được.

Từ $A_S$ bán chuẩn hóa
(chuẩn hóa) tối đa các cột
\[H = \begin{pmatrix} \gamma_1 \\ \gamma_2 \\ \vdots \\ \gamma_k \\ \text{O} \\ \vdots \\ \text{O} \end{pmatrix}\]

\emph{k} dòng khác zero

các dòng zero, có thể có hoặc không

Suy ra tập hợp $a = \{\gamma_1, \gamma_2, \dots, \gamma_k\}$ là một
cơ sở của \emph{W} và số chiều là
$\dim_F W = k$.

\ul{Ví dụ}: $A_S$bán chuẩn
hóa (chuẩn hóa) tối đa các cột
\[H = \begin{pmatrix} -1 & 0 & 2 & -3 & 4 & 0 & 1 \\ 0 & 2 & -5 & 1 & 1 & -3 & 2 \\ 0 & 0 & 0 & -2 & 1 & 0 & 8 \\ 0 & 0 & 0 & 0 & 0 & 0 & 0 \\ 0 & 0 & 0 & 0 & 0 & 0 & 0 \end{pmatrix} = \begin{pmatrix} \gamma_1 \\ \gamma_2 \\ \gamma_3 \\ \text{O} \\ \text{O} \end{pmatrix}\]

Suy ra
\[ a = \{\gamma_1 = (-1; 0; 2; -3; 4; 0; 1), \gamma_2 = (0; 2; -5; 1; 1; -3; 2), \gamma_3 = (0; 0; 0; -2; 1; 0; 8)\} \]
là một cơ sở
của \emph{W}, và số chiều là
$\dim_{\mathbb{R}} W = 3$.

\ul{Bài tập}: Tìm cơ sở và xác định số chiều cho không gian \emph{W}
trong các trường hợp sau:

a/ $V = \mathbb{R}^3$,
\[ S = \{a_1 = (1; -3; 2), a_2 = (0; 1; -4), a_3 = (-2; 3; 1)\} \]
và
$W = \langle S \rangle$.

b/ $V = \mathbb{R}^4$,
\[ S = \{a_1 = (1; 0; -1; 2), a_2 = (2; 0; 3; -1), a_3 = (-1; -2; 0; 2)\} \]
và
$W = \langle S \rangle$.

c/ $V = \mathbb{R}^4$,
\[ S = \{a_1 = (-1; 3; 2; -4), a_2 = (0; 0; 1; -1), a_3 = (1; 0; 1; 0), a_4 = (2; 3; -1; 4)\} \]
và
$W = \langle S \rangle$.

d/ $V = \mathbb{R}^5$,
\[ S = \{a_1 = (-1; 1; 0; -1; 1), a_2 = (1; 0; 0; 2; -1), a_3 = (-2; 1; 3; 0; -1)\} \]
và
$W = \langle S \rangle$.

e/ $V = \mathbb{R}^5$,
\[ S = \{a_1 = (-2; -3; 0; 0; 1), a_2 = (-1; 1; 1; 0; 0), a_3 = (2; -1; 4; 1; 0), a_4 = (-2; 1; -3; 0; 1)\} \]
và
$W = \langle S \rangle$

\textbf{* \ul{TÌM HỆ SINH, CƠ SỞ VÀ XÁC ĐỊNH SỐ CHIỀU CHO KHÔNG GIAN PHỤ
THUỘC THAM SỐ}}.

\ul{Ví dụ mẫu}: Trên $V = \mathbb{R}^4$,
cho $W = \{\alpha = (a - b + 2c; 2b - a - c; 3c + 2a; c + a - 3b) \mid a, b, c \in \mathbb{R}\}$

a/ Chứng minh rằng $W \le V$.

b/ Tìm hệ sinh, cơ sở và xác định số chiều cho \emph{W}.

\ul{Giải}:

Ta có $W = \{\alpha = (a - b + 2c; 2b - a - c; 3c + 2a; c + a - 3b) \mid a, b, c \in \mathbb{R}\}$

$= \{\alpha = (a; -a; 2a; a) + (-b; 2b; 0b; -3b) + (2c; -c; 3c; c) \mid a, b, c \in \mathbb{R}\}$

$= \{\alpha = a(1; -1; 2; 1) + b(-1; 2; 0; -3) + c(2; -1; 3; 1) \mid a, b, c \in \mathbb{R}\}$

Đặt
\[ \alpha_1 = (1; -1; 2; 1), \alpha_2 = (-1; 2; 0; -3), \alpha_3 = (2; -1; 3; 1) \]
và gọi $S = \{\alpha_1, \alpha_2, \alpha_3\}$.

Suy ra $W = \{ \alpha = a\alpha_1 + b\alpha_2 + c\alpha_3 \mid a, b, c \in \mathbb{R}; \alpha_1, \alpha_2, \alpha_3 \in S \}$

$= \langle S \rangle$

Mà $\langle S \rangle \le V = \mathbb{R}^4$ nên ta nói
$W = \langle S \rangle \le V$.

Ta có $W = \langle S \rangle$ nên
$S = \{\alpha_1, \alpha_2, \alpha_3\}$ là một hệ sinh của
\emph{W}.

Tiếp theo tìm cơ sở cho \emph{W} như sau:

Ta lập ma trận

\begin{quote}
\[A_S = \begin{pmatrix} \alpha_1 \\ \alpha_2 \\ \alpha_3 \end{pmatrix} = \begin{pmatrix} 1 & -1 & 2 & 1 \\ -1 & 2 & 0 & -3 \\ 2 & -1 & 3 & 1 \end{pmatrix} \xrightarrow{\substack{(2) \to (2) + (1) \\ (3) \to (3) - 2(1)}} \begin{pmatrix} 1 & -1 & 2 & 1 \\ 0 & 1 & 2 & -2 \\ 0 & 1 & -1 & -1 \end{pmatrix}\]
\end{quote}

\[\xrightarrow{(3) \to (3) - (2)} \begin{pmatrix} 1 & -1 & 2 & 1 \\ 0 & 1 & 2 & -2 \\ 0 & 0 & -3 & 1 \end{pmatrix} = \begin{pmatrix} \gamma_1 \\ \gamma_2 \\ \gamma_3 \end{pmatrix}\]

\ul{Cách 1}: Do ma trận kết quả có 3 dòng khác zero
$r(A_S) = 3 =$số véc tơ trong
\emph{S},

Cho nên \emph{S} là một hệ sinh, ĐLTT, nên \emph{S} cũng là một cơ sở
của \emph{W}, và số chiều là

$\dim_{\mathbb{R}} W = 3$.

\ul{Cách 2}: Ta có $a = \{ y_1 = (1; -1; 2; 1), y_2 = (0; 1; 2; -2), y_3 = (0; 0; -3; 1) \}$ là
một cơ sở của \emph{W},

\begin{quote}
và số chiều là
\end{quote}

$\dim_{\mathbb{R}} W = 3$.

\emph{\ul{Bài tập tương tự}}: Yêu cầu giống như ví dụ mẫu phần tìm hệ
sinh, cơ sở và xác định số chiều cho không gian phụ thuộc tham số:

a/ $V = \mathbb{R}^4$,
\[ W = \{ \alpha = (-a + 2b + c; 4c + 5a; 2c - b + a; 2a - 3b - c) \mid a, b, c \in \mathbb{R} \}. \]

b/ $V = \mathbb{R}^4$,
\[ W = \{ \alpha = (3d + a - 2b - c; 2d - c + a; 3d - c + b + 2a; d - a + b) \mid a, b, c, d \in \mathbb{R} \}. \]

c/ $V = \mathbb{R}^4$,
\[ W = \{ \alpha = (d - e + 2a + b - 3c; 4e + d - 2c + a; e - 3d + c + 2b - a; 3e + 2d - a - b) \mid a, b, c, d, e \in \mathbb{R} \}. \]

d/ $V = \mathbb{R}^5$,
\[ W = \{ \alpha = (a - 2b + c; 2c - 3a + b; 4c + b - 2a; a - b + 3c; c - 2a - b) \mid a, b, c \in \mathbb{R} \}. \]

e/ $V = \mathbb{R}^5$,
\[ W = \{ \alpha = (a - b + 2c - d; c + a - 2b + 3d; 2c + d - a; a - 2b + c - 3d; c + d - a + b) \mid a, b, c, d \in \mathbb{R} \}. \]

\textbf{* \ul{TÌM HỆ SINH, CƠ SỞ VÀ XÁC ĐỊNH SỐ CHIỀU CHO KHÔNG GIAN
NGHIỆM CỦA HỆ PHƯƠNG TRÌNH TUYẾN TÍNH THUẦN NHẤT AX = O}}.

Cho $(V, +, \bullet)$ là không gian
tuyến tính tổng quát trên trường số
$F = \mathbb{R}$,

và gọi $W = \{ X \in \mathbb{R}^n \mid AX = O \}$.

Ta sẽ chứng minh $W \le V$ và
tìm hệ sinh, cơ sở, số chiều cho
$W$ như sau:

+ Từ hệ phương trình thuần nhất
$AX = O$, ta giải hệ pt này
(nên dùng pp Gauss hoặc Gauss-Jordan), để tìm nghiệm.

+ Ta thay nghiệm tìm được vào
$W$, và viết lại
$W$ dưới dạng không gian
phụ thuộc tham số.

+ Sau đó ta chứng tỏ $W \le V$
và tìm hệ sinh, cơ sở, số chiều cho
$W$ hoàn toàn giống như
trường hợp tìm hệ sinh, cơ sở và xác định số chiều cho không gian phụ
thuộc tham số đã nêu trên.

\ul{Ví dụ mẫu 1}: Trên không gian tuyến tính
$V = \mathbb{R}^5$, cho

\[W = \left\{ X = (x_1; x_2; x_3; x_4; x_5) \in \mathbb{R}^5 \left| \begin{matrix} 2x_3 + 7x_5 - x_4 + x_1 + 3x_2 = 0 \\ x_2 + 3x_4 - x_5 + 2x_1 - x_3 = 0 \\ 5x_3 - 8x_4 - 3x_1 + 2x_2 + 12x_5 = 0 \\ 7x_4 - 9x_5 + 3x_1 - x_2 - 4x_3 = 0 \end{matrix} \right. \right\}\]

a/ Chứng minh rằng $W \le V$.

b/ Tìm hệ sinh, cơ sở và xác định số chiều cho
$W$.

\ul{Giải}:

Từ điều kiện của \emph{W} ta viết lại hệ pt tuyến tính thuần nhất dưới
dạng ma trận hóa như sau:

\begin{align*}
&\begin{pmatrix} 1 & 3 & 2 & -1 & 7 \mid 0 \\ 2 & 1 & -1 & 3 & -1 \mid 0 \\ -3 & 2 & 5 & -8 & 12 \mid 0 \\ 3 & -1 & -4 & 7 & -9 \mid 0 \end{pmatrix} \\
&\xrightarrow{\substack{(2) \to (2) - 2(1) \\ (3) \to (3) + 3(1) \\ (4) \to (4) - 3(1)}}
\begin{pmatrix} 1 & 3 & 2 & -1 & 7 \mid 0 \\ 0 & -5 & -5 & 5 & -15 \mid 0 \\ 0 & 11 & 11 & -11 & 33 \mid 0 \\ 0 & -10 & -10 & 10 & -30 \mid 0 \end{pmatrix} \\
&\xrightarrow{\substack{(2) \to -\frac{1}{5}(2) \\ (3) \to (3) - 11(2)_{moi} \\ (4) \to (4) + 10(2)_{moi}}}
\begin{pmatrix} 1 & 3 & 2 & -1 & 7 \mid 0 \\ 0 & 1 & 1 & -1 & 3 \mid 0 \\ 0 & 0 & 0 & 0 & 0 \mid 0 \\ 0 & 0 & 0 & 0 & 0 \mid 0 \end{pmatrix}
\end{align*}

Do cột 3,4,5, không bán chuẩn hóa được, nên hệ pt có vô số nghiệm như
sau:

Đặt \[\begin{cases} x_1 = a \\ x_4 = b, \forall a, b, c \in \mathbb{R} \\ x_5 = c \end{cases}\].

Ta có hệ pt lúc này:

\begin{align*}
\begin{cases} x_1 + 3x_2 + 2x_3 - x_4 + 7x_5 = 0 \\ x_2 + x_3 - x_4 + 3x_5 = 0 \end{cases} 
&\Rightarrow \begin{cases} x_1 = -3(-a+b-3c) - 2a + b - 7c \\ x_2 = -a + b - 3c \end{cases} \\
&\Rightarrow \begin{cases} x_1 = a - 2b + 2c \\ x_2 = -a + b - 3c \end{cases}
\end{align*}

Thay vào \emph{W} ta được:

$W = \{ X = (x_1; x_2; x_3; x_4; x_5) = (a - 2b + 2c; -a + b - 3c; a; b; c) \mid a, b, c \in \mathbb{R} \}$

$= \{ X = (a; -a; a; 0a; 0a) + (-2b; b; 0b; b; 0b) + (2c; -3c; 0c; 0c; c) \mid a, b, c \in \mathbb{R} \}$

$= \{ X = a(1; -1; 1; 0; 0) + b(-2; 1; 0; 1; 0) + c(2; -3; 0; 0; 1) \mid a, b, c \in \mathbb{R} \}$.

Đặt
\[ \alpha_1 = (1; -1; 1; 0; 0), \alpha_2 = (-2; 1; 0; 1; 0), \alpha_3 = (2; -3; 0; 0; 1) \]
và gọi $S = \{ \alpha_1, \alpha_2, \alpha_3 \}$.

$\Rightarrow W = \{ X = a\alpha_1 + b\alpha_2 + c\alpha_3 \mid a, b, c \in \mathbb{R}; \alpha_1, \alpha_2, \alpha_3 \in S \}$

$= \langle S \rangle$.

Mà $\langle S \rangle \le V$ nên
$W = \langle S \rangle \le V$.

Ta có $W = \langle S \rangle$ nên
$S = \{ \alpha_1, \alpha_2, \alpha_3 \}$ là một hệ sinh của
\emph{W}.

Tiếp theo, ta tìm cơ sở và số chiều cho \emph{W} như sau:

Lập ma trận:

\ul{Cách 1}: Do ma trận đáp số có 3 dòng khác zero nên ta nói
$r(A_S) = 3 =$số véc tơ trong
\emph{S}, nên ta nói \emph{S} vừa là hệ sinh, vừa ĐLTT nên \emph{S} cũng
là 1 cơ sở của \emph{W}, và số chiều là
$\dim_{\mathbb{R}} W = 3$.

\ul{Cách 2}: Ta có $a = \{ \gamma_1 = (1; -1; 1; 0; 0), \gamma_2 = (0; -2; 0; 1; 1), \gamma_3 = (0; 0; -4; -1; 1) \}$ là
1 cơ sở của \emph{W}, và số chiều là
$\dim_{\mathbb{R}} W = 3$.

\emph{\ul{Bài tập tương tự}}: Yêu cầu như Ví dụ mẫu 1

a/ \[V = \mathbb{R}^4, W = \left\{ X = (x_1; x_2; x_3; x_4) \left| \begin{matrix} 2x_2 - x_1 + x_3 = 0 \\ x_1 - 3x_2 + 2x_4 = 0 \\ 2x_1 - 3x_2 - x_4 = 0 \end{matrix} \right. \right\}\].

b/ \[V = \mathbb{R}^4, W = \left\{ X = (x_1; x_2; x_3; x_4) \left| \begin{matrix} x_2 + x_1 + 2x_3 = 0 \\ 2x_2 - x_1 - x_4 = 0 \end{matrix} \right. \right\}\].

c/ \[V = \mathbb{R}^5, W = \left\{ X = (x_1; x_2; x_3; x_4; x_5) \left| \begin{matrix} x_5 - x_1 + x_2 + 2x_3 = 0 \\ 4x_5 - 2x_1 - x_2 + 2x_3 = 0 \\ x_1 - 3x_5 + x_2 - 2x_4 = 0 \end{matrix} \right. \right\}\].

d/ \[V = \mathbb{R}^5, W = \left\{ X = (x_1; x_2; x_3; x_4; x_5) \left| \begin{matrix} x_5 - x_4 - x_1 + 2x_2 = 0 \\ 3x_5 - x_1 + 4x_2 - x_3 = 0 \\ x_1 + x_2 - 2x_5 = 0 \end{matrix} \right. \right\}\].

e/ \[V = \mathbb{R}^5, W = \left\{ X = (x_1; x_2; x_3; x_4; x_5) \left| \begin{matrix} 2x_2 - x_5 + x_1 - x_3 = 0 \\ x_4 + x_1 + 3x_5 = 0 \\ 2x_1 - x_2 - 3x_5 + x_4 = 0 \\ 2x_3 - x_2 - 6x_5 + x_1 = 0 \end{matrix} \right. \right\}\].

\textbf{7/ \ul{TỌA ĐỘ VÉC TƠ THEO CƠ SỞ}}:

Cho $(V, +, \bullet)$ là không gian
tuyến tính tổng quát trên trường số \emph{F}.

Gọi $S = \{ \alpha_1, \alpha_2, \dots, \alpha_m \}$ là một cơ sở của
\emph{V}.

Cho trước véc tơ $\alpha \in V$. Lúc
này, luôn tồn tại bộ số
$c_1, c_2, \dots, c_m \in F$ sao cho:

$\alpha = c_1\alpha_1 + c_2\alpha_2 + \dots + c_m\alpha_m$

(ở đây các ẩn số là $c_1, c_2, \dots, c_m$;
và ta cần giải hệ pt để tìm bộ nghiệm này).

Khi đó ta viết:

\[[\alpha]_S = \begin{bmatrix} c_1 \\ c_2 \\ \vdots \\ c_m \end{bmatrix} = \begin{pmatrix} c_1 \\ c_2 \\ \vdots \\ c_m \end{pmatrix} =\] đây là tọa độ của véc
tơ $\alpha$ theo cơ sở
\emph{S}.

\ul{Ví dụ mẫu 2}: Trên
$V = \mathbb{R}^3$, cho hệ
$S = \{ a_1 = (-2; 4; 3), a_2 = (-4; 5; 1), a_3 = (1; 1; 2) \}$.

a/ Chứng minh rằng \emph{S} là một cơ sở của
$\mathbb{R}^3$.

b/ Cho véc tơ $\alpha = (7; -8; 1) \in \mathbb{R}^3$. Tìm
$[\alpha]_s =?$

c/ Cho véc tơ $\beta \in \mathbb{R}^3$ thỏa
\[[\beta]_s = \begin{bmatrix} -4 \\ 1 \\ 3 \end{bmatrix}\]. Tìm
$\beta =?$

d/ Cho véc tơ $\gamma \in \mathbb{R}^3$ thỏa
\[[\gamma]_s = \begin{bmatrix} 2 \\ -5 \\ 7 \end{bmatrix}\]. Tìm
$\gamma =?$

\ul{Giải}:

a/ Do $\mathbb{R}^3$ là không gian
tuyến tính 3 chiều trên trường số thực
$\mathbb{R}$, nên mỗi cơ sở của
$\mathbb{R}^3$ là một hệ sinh, có 3
véc tơ, ĐLTT. Mà tập hợp \emph{S} đã có 3 véc tơ nên \emph{S} là 1 hệ
sinh của $\mathbb{R}^3$. Cho nên để
chứng tỏ \emph{S} là một cơ sở của
$\mathbb{R}^3$ thì ta chỉ cần chứng
tỏ \emph{S} là ĐLTT.

Ta lập ma trận

\begin{quote}
\[A_S = \begin{pmatrix} a_1 \\ a_2 \\ a_3 \end{pmatrix} = \begin{pmatrix} -2 & 4 & 3 \\ -4 & 5 & 1 \\ 1 & 1 & 2 \end{pmatrix}\]
\end{quote}

Ta có

\begin{quote}
$\det(A_S) = (-2).5.2 + 4.1.1 + 3.(-4).1 - [3.5.1 + 4.(-4).2 + (-2).1.1] = -9 \ne 0$
\end{quote}

$\Rightarrow S$ là một hệ sinh ĐLTT,
nên \emph{S} là một cơ sở của
$\mathbb{R}^3$.

b/ Cho véc tơ $\alpha = (7; -8; 1) \in \mathbb{R}^3$. Tìm
$[\alpha]_s =?$

Ta cần tìm $c_1, c_2, c_3 \in \mathbb{R}$ sao cho:
$\alpha = c_1a_1 + c_2a_2 + c_3a_3$

\begin{align*}
(7, -8, 1) &= c_1(-2, 4, 3) + c_2(-4, 5, 1) + c_3(1, 1, 2) \\
&\Leftrightarrow \begin{cases} -2c_1 - 4c_2 + c_3 = 7 \\ 4c_1 + 5c_2 + c_3 = -8 \\ 3c_1 + c_2 + 2c_3 = 1 \end{cases} \\
&\Leftrightarrow \begin{cases} c_1 = 2 \\ c_2 = -3 \\ c_3 = -1 \end{cases}
\end{align*}

Như vậy, ta có:

\begin{quote}
\[[\alpha]_S = \begin{pmatrix} c_1 \\ c_2 \\ c_3 \end{pmatrix} = \begin{pmatrix} 2 \\ -3 \\ -1 \end{pmatrix}\]
\end{quote}

c/ Cho véc tơ $\beta \in \mathbb{R}^3$ thỏa
\[[\beta]_s = \begin{bmatrix} -4 \\ 1 \\ 3 \end{bmatrix}\]. Tìm
$\beta =?$

Tìm $\beta = -4a_1 + 1a_2 + 3a_3 = -4(-2; 4; 3) + (-4; 5; 1) + 3(1; 1; 2) = (7; -8; -5)$.

d/ Cho véc tơ $\gamma \in \mathbb{R}^3$ thỏa
\[[\gamma]_s = \begin{bmatrix} 2 \\ -5 \\ 7 \end{bmatrix}\]. Tìm
$\gamma =?$.

Tìm $\gamma = 2a_1 - 5a_2 + 7a_3 = 2(-2; 4; 3) - 5(-4; 5; 1) + 7(1; 1; 2) = (23; -10; 15)$.

* \emph{\textbf{\ul{Lưu ý}}}: Trên
$V = \mathbb{R}^n$, luôn có cơ sở chính
tắc

\begin{quote}
$\beta_0 = \{e_1 = (1, 0, \dots, 0), e_2 = (0, 1, 0, \dots, 0), \dots, e_n = (0, 0, \dots, 1)\}$
\end{quote}

Với mọi véc tơ $\alpha = (x_1; x_2; ...; x_n) \in \mathbb{R}^n$ thì ta
có

\[[\alpha]_{\beta_0} = \begin{bmatrix} x_1 \\ x_2 \\ \vdots \\ x_n \end{bmatrix}\]

(lấy tọa độ của $\alpha$ viết
thành 1 cột mà ta không cần phải giải hệ pt để tìm các số

$c_1, c_2, ..., c_n$).

\ul{Ví dụ}: Trên $V = \mathbb{R}^3$, ta
luôn có cơ sở chính tắc là

$\beta_0 = \{e_1 = (1, 0, 0), e_2 = (0, 1, 0), e_3 = (0, 0, 1)\}$

Thì với $\alpha = (7; -8; 1) \in \mathbb{R}^3$, ta có
\[[\alpha]_{\beta_0} = \begin{bmatrix} 7 \\ -8 \\ 1 \end{bmatrix}\]

$\lambda = (0; 2; -5) \in \mathbb{R}^3$, ta có
\[[\lambda]_{\beta_0} = \begin{bmatrix} 0 \\ 2 \\ -5 \end{bmatrix}\]

\ul{Ví dụ mẫu 3}: bài 3.18

Bài 3.18. Hãy tìm tọa độ của véc tơ $x = (8, 8, 19, 19)$
trong cơ sở dưới đây của không gian tuyến tính $\mathbb{R}^4$:
$a_1 = (1, 1, 2, 3);\; a_2 = (2, 1, 3, 4);$
$a_3 = (2, 3, -2, 1);\; a_4 = (1, 3, 3, 1).$

\ul{Giải}:

Do cơ sở của $\mathbb{R}^4$ có 4 véc
tơ $a_1, a_2, a_3, a_4$ nên ta cần tìm 4
số $c_1, c_2, c_3, c_4 \in \mathbb{R}$ sao cho

$x = c_1a_1 + c_2a_2 + c_3a_3 + c_4a_4$

$\Leftrightarrow (8, 8, 19, 19) = c_1(1, 1, 2, 3) + c_2(2, 1, 3, 4) + c_3(2, 3, -2, 1) + c_4(1, 3, 3, 1)$

\[\Leftrightarrow \begin{cases} c_1 + 2c_2 + 2c_3 + c_4 = 8 \\ c_1 + c_2 + 3c_3 + 3c_4 = 8 \\ 2c_1 + 3c_2 - 2c_3 + 3c_4 = 19 \\ 3c_1 + 4c_2 + c_3 + c_4 = 19 \end{cases} \Leftrightarrow \begin{cases} c_1 = \frac{92}{35} \\ c_2 = \frac{89}{35} \\ c_3 = -\frac{23}{35} \\ c_4 = \frac{8}{5} \end{cases}\]

Ta có:

\begin{quote}
\[[x]_S = \begin{pmatrix} \frac{92}{35} \\ \frac{89}{35} \\ -\frac{23}{35} \\ \frac{8}{5} \end{pmatrix}\]
\end{quote}

\ul{Ví dụ mẫu 4}: bài 3.20

Bài 3.20. Trong không gian tuyến tính $\mathbb{R}^4$ cho $M$ là
không gian con hai chiều có cơ sở là $\{u_1, u_2\}$ với
$u_1 = (2, 1, -1, 1),\quad u_2 = (1, 2, 3, -1).$
Cho các phần tử $u = (0, 1, 1, 3),\; v = (1, 1, 1, -1)$. Hãy
xác định số thực $\lambda$ sao cho $u - \lambda v \in M.$

\ul{Giải}:

Gọi $S = \{u_1 = (2; 1; -1; 1), u_2 = (1; 2; 3; -1)\}$ là cơ sở của
không gian $M = \langle S \rangle$

Tìm số thực $\lambda$ sao cho
$(u - \lambda v) \in M$.

Yêu cầu đề bài tương đương tìm số thực
$\lambda$ sao cho

$u - \lambda v = c_1u_1 + c_2u_2$ có nghiệm

$\Leftrightarrow (0, 1, 1, 3) - \lambda(1, 1, 1, -1) = c_1(2, 1, -1, 1) + c_2(1, 2, 3, -1)$

$\Leftrightarrow (-\lambda, 1 - \lambda, 1 - \lambda, 3 + \lambda) = (2c_1 + c_2, c_1 + 2c_2, -c_1 + 3c_2, c_1 - c_2)$

\[\Leftrightarrow \begin{cases} 2c_1 + c_2 = -\lambda \\ c_1 + 2c_2 = 1 - \lambda \\ -c_1 + 3c_2 = 1 - \lambda \\ c_1 - c_2 = 3 + \lambda \end{cases}\]

Viết lại hệ pt dưới dạng ma trận hóa ta được:

\begin{align*}
&\begin{pmatrix} 2 & 1 \mid -\lambda \\ 1 & 2 \mid 1 - \lambda \\ -1 & 3 \mid 1 - \lambda \\ 1 & -1 \mid 3 + \lambda \end{pmatrix} \\
&\xrightarrow{\substack{(2) \to 2(2) - (1) \\ (3) \to 2(3) + (1) \\ (4) \to 2(4) - (1)}} \begin{pmatrix} 2 & 1 \mid -\lambda \\ 0 & 3 \mid 2 - 2\lambda \\ 0 & 7 \mid 2 - 3\lambda \\ 0 & -3 \mid 6 + 3\lambda \end{pmatrix} \\
&\xrightarrow{\substack{(3) \to 3(3) - 7(2) \\ (4) \to (4) + (2)}} \begin{pmatrix} 2 & 1 \mid -\lambda \\ 0 & 3 \mid 2 - 2\lambda \\ 0 & 0 \mid -8 - \lambda \\ 0 & 0 \mid 8 + \lambda \end{pmatrix}
\end{align*}

Để hệ có nghiệm thì dòng (3) và dòng (4) phải thỏa dạng
\[\begin{pmatrix} 0 & 0 | 0 \\ 0 & 0 | 0 \end{pmatrix}\] nghĩa là ta có

\[\begin{cases} 16 + 4\lambda = 0 \\ 36 + 9\lambda = 0 \end{cases} \Leftrightarrow \lambda = -4\]

Đáp số: $\lambda = -4$ là giá trị
cần tìm.

\ul{Ví dụ mẫu 5}: bài 3.24 trang 5

Bài 3.24. Trong không gian tuyến tính $\mathbb{R}^4$ cho không
gian con
\[ M = \{(x_1, x_2, x_3, x_4) \mid 2x_1 - x_2 - x_3 + 4x_4 = 0\} \]
và phần tử $w \in M$ với $w = (1, 5, 1, 1)$. Hãy xác định
một cơ sở và số chiều của $M$ và cho biết tọa độ của
$w$ trên cơ sở được đưa ra.

\ul{Giải}:

Từ điều kiện của \emph{M} ta có hệ pt thuần nhất:

$2x_1 - x_2 - x_3 + 4x_4 = 0$

Viết hệ dưới dạng ma trận hóa ta được:

$(2 \quad -1 \quad -1 \quad 4 \mid 0)$

Do cột 2, 3, 4 không bán chuẩn hóa được nên hệ có vô số nghiệm như sau:

Đặt \[\begin{cases} x_2 = 2a \\ x_3 = 2b, \forall a,b,c \in \mathbb{R} \\ x_4 = c \end{cases}\].

Ta có $2x_1 - x_2 - x_3 + 4x_4 = 0 \Rightarrow x_1 = a + b - 2c$

$\Rightarrow M = \{(x_1; x_2; x_3; x_4) = (a + b - 2c; 2a; 2b; c) | a,b,c \in \mathbb{R}\}$

$= \{ \alpha = (a; 2a; 0a; 0a) + (b; 0b; 2b; 0b) + (-2c; 0c; 0c; c) | a,b,c \in \mathbb{R} \}$

$= \{ \alpha = a(1; 2; 0; 0) + b(1; 0; 2; 0) + c(-2; 0; 0; 1) | a,b,c \in \mathbb{R} \}$

Đặt
\[ a_1 = (1; 2; 0; 0), a_2 = (1; 0; 2; 0), a_3 = (-2; 0; 0; 1) \]
và gọi $S = \{ \alpha_1, \alpha_2, \alpha_3 \}$

$\Rightarrow M = \{ \alpha = a\alpha_1 + b\alpha_2 + c\alpha_3 \mid a, b, c \in \mathbb{R}; \alpha_1, \alpha_2, \alpha_3 \in S \}$

$= \langle S \rangle$

Cho nên $M = \langle S \rangle$ nên
$S = \{ \alpha_1, \alpha_2, \alpha_3 \}$ là một hệ sinh của
\emph{M}.

Tiếp theo, ta tìm cơ sở và số chiều cho \emph{M} như sau:

Lập ma trận

\[A_S = \begin{pmatrix} \alpha_1 \\ \alpha_2 \\ \alpha_3 \end{pmatrix} = \begin{pmatrix} 1 & 2 & 0 & 0 \\ 1 & 0 & 2 & 0 \\ -2 & 0 & 0 & 1 \end{pmatrix} \xrightarrow{\substack{(2) \to (2) - (1) \\ (3) \to (3) + 2(1)}} \begin{pmatrix} 1 & 2 & 0 & 0 \\ 0 & -2 & 2 & 0 \\ 0 & 4 & 0 & 1 \end{pmatrix} = \begin{pmatrix} \gamma_1 \\ \gamma_2 \\ \gamma_3 \end{pmatrix}\]

Ta có: $a = \{ \gamma_1 = (1; 2; 0; 0), \gamma_2 = (0; -2; 2; 0), \gamma_3 = (0; 0; 4; 1) \}$ là một cơ sở
của \emph{M}, và số chiều là:

$\dim_{\mathbb{R}} M = 3$.

Tìm tọa độ của $w = (1; 5; 1; 1)$ theo
cơ sở \emph{a} của \emph{M} vừa tìm được.

Ta tìm $c_1, c_2, c_3 \in \mathbb{R}$ sao cho
$w = c_1 \gamma_1 + c_2 \gamma_2 + c_3 \gamma_3$

$\Leftrightarrow (1, 5, 1, 1) = c_1(1, 2, 0, 0) + c_2(0, -2, 2, 0) + c_3(0, 0, 4, 1)$

\[\Leftrightarrow \begin{cases} c_1 = 1 \\ 2c_1 - 2c_2 = 5 \\ 2c_2 + 4c_3 = 1 \\ c_3 = 1 \end{cases} \Rightarrow \begin{cases} c_1 = 1 \\ c_2 = -\frac{3}{2} \\ c_3 = 1 \end{cases}\]

Vậy, \[[w]_a = \begin{bmatrix} c_1 \\ c_2 \\ c_3 \end{bmatrix} = \begin{bmatrix} 1 \\ -3/2 \\ 1 \end{bmatrix}\].

\end{document}
